\section{The constants produced by the proof of the polynomial PEPS approximation}
\label{sec:polynomial_peps_explicit_constants}

Theorem~1.1 of \cite{OpenAI2026PolynomialPEPS} asserts that its constants exist, and Section~8 fixes them in a definite order (\texttt{07-assembly.tex}, lines 179--214). Below we express the large-size constants as explicit functions of the square area-law constant $C_A$ and the angular slope $D$ and threshold $L_{\rm ang}$ of Proposition~3.4.

% Status: these are auxiliary results. The source asserts only that the
% constants exist; the statement of Theorem 1.1 is thm:openai_polynomial_peps.

% Source: 06-geometry.tex, lines 12--53 and 559--681.
% Revision: adc7f1241b42e322a6451854ab7e4b4c146bf78a.
\begin{lemma}[The distribution protocol at a given slope and given patch constants]
    \label{lem:peps_dyadic_protocol_at_given_scale}
    \lean{TNLean.PEPS.Approximation.largeSizeDyadicProtocol_of_angularScale}
    \leanok
    \uses{def:peps_dyadic_protocol,def:peps_angular_scale,
        lem:peps_adaptive_patch_supply,
        lem:peps_dyadic_protocol_mesh_ratio}
    Fix $q\ge2$ and a family of unit states on the $L\times L$ squares. Set $\epsilon_0=1/30000$, $\nu=1/600000000$, $a_0=1/5000$, and $b=\max\{18,\lceil((4+8\epsilon_0+2\nu)/\nu+1)^2\rceil\}$. Let $D\ge1$ and $L_{\rm ang}$ be an angular slope and threshold for the family at accuracy $L^{-60}$, clearance $a_0/2$, diameter ratio $1$, and at most $4$ marked points, and let $C>0,c\in\mathbb R$ be adaptive-patch constants (Proposition~4.1) on the cell-center sites at radius constant $B_0=20D$ and power $60$ for all $L\ge2$. With $C'=\lceil C\rceil$ and $c'=\lceil c\rceil$, set
    \begin{align}
        G   & =(C'+1)^{3b}\bigl(16b^2(C'+1)^{6b}+1\bigr)^{2b}, &
        k_K & =3bc'+4b(3bc'+30).
    \end{align}
    Then for every $C_N\ge5042$, $C_K\ge G$, $C_E\ge\max\{10080,2b+1\}$, $L\ge\max\{2,\lceil\max\{L_{\rm ang},1\}\rceil\}$, and family state $\Omega$ at side $L$, there is a protocol satisfying Definition~\ref{def:peps_dyadic_protocol} with constants $(C_N,C_K,C_E)$, party bound $64\cdot189$, exponent $k_K$, and $2b$ pair effects.
\end{lemma}

\begin{proof}\leanok
    \uses{def:peps_dyadic_protocol_parties,
        def:peps_protocol_frozen_frames,
        thm:peps_angular_scale_monotonicity,
        thm:peps_dyadic_direct_repainting_constants,
        thm:peps_dyadic_padding_bounds,
        thm:peps_dyadic_protocol_schedule_counts,
        thm:peps_protocol_constants,
        thm:peps_protocol_from_large_repaintings,
        thm:peps_protocol_frozen_frame_bounds,
        thm:peps_protocol_large_repainting_segment,
        thm:peps_protocol_level_events,
        thm:peps_protocol_point_treatment_change}
    This follows the proof of Theorem~\ref{thm:peps_protocol_from_angular_and_patch_bounds} using the supplied $(D,L_{\rm ang},C,c)$: the mesh ratio $\nu$ satisfies the resize, point-treatment, and direct-repainting conditions for $K_0=20$, $a=a_0/2$, and $B_0=K_0D$, and $L\ge\lceil\max\{L_{\rm ang},1\}\rceil$ activates the angular slope at side $L$.
\end{proof}

% Source: 07-assembly.tex, lines 45--63, 122--177 and 197--206.
% Revision: adc7f1241b42e322a6451854ab7e4b4c146bf78a.
\begin{definition}[The constants of the polynomial PEPS approximation for large sizes]
    \label{def:polynomial_peps_explicit_constants}
    \lean{TNLean.PEPS.Approximation.pepsLinkExponentOf,
        TNLean.PEPS.Approximation.pepsLinkPrefactorOf,
        TNLean.PEPS.Approximation.pepsCongestion,
        TNLean.PEPS.Approximation.pepsCongestion_eq,
        TNLean.PEPS.Approximation.pepsNetworkErrorConstant,
        TNLean.PEPS.Approximation.pepsSizeThresholdOf,
        TNLean.PEPS.Approximation.pepsBondExponentOf,
        TNLean.PEPS.Approximation.pepsLargeSizePrefactorOf,
        TNLean.PEPS.Approximation.pepsLargeSizePrefactorOf_pos,
        TNLean.PEPS.Approximation.pepsBondExponentOf_pos,
        TNLean.PEPS.Approximation.pepsLinkExponentOfPatch,
        TNLean.PEPS.Approximation.pepsBondExponentOfPatch,
        TNLean.PEPS.Approximation.pepsBondExponentOfPatch_pos,
        TNLean.PEPS.Approximation.pepsBondExponentOf_eq_ofPatch}
    \leanok
    \uses{def:peps_adaptive_patch_constants,def:peps_compression_constants,
        def:peps_sampled_network_constants,
        thm:peps_dyadic_protocol_explicit_constants}
    Fix $q\in\Z$ and $J,\Delta,C_A,D,L_{\rm ang}\in\R$. Set $\epsilon_0=1/30000$, $\nu=1/600000000$, $b_{\rm prot}=64\cdot189$, and
    \begin{align}
        b=\max\{18,\lceil((4+8\epsilon_0+2\nu)/\nu+1)^2\rceil\}
        =5760768040000960009.
    \end{align}
    Let $C=\max\{q^{81},1+10D\}$ and $c=c(q,J,\Delta,20D,60,C_A)$ be the constants of Definition~\ref{def:peps_adaptive_patch_constants} at $B_0=20D$, $p=60$, and $C_A$; let $G,k_K$ be as in Theorem~\ref{thm:peps_dyadic_protocol_explicit_constants} for $C'=\lceil C\rceil$, $c'=\lceil c\rceil$; and set $C_E=\max\{10080,2b+1\}$. Define the link exponent $k$ and prefactor $C_1$ (from Definitions~\ref{def:peps_sampled_network_constants} and~\ref{def:peps_compression_constants} at $2b$ pair effects, arity $b_{\rm prot}$, gate-count exponent $3$, physical-dimension exponent $0$, monomial/coefficient-sum exponent $k_K$, and accuracy $L^{-5}$), congestion $\chi$, error constant $C_{\rm net}$, size threshold $L_0$, and large-size bond constants $c_{\rm PEPS},C_{\rm large}>0$ by
    \begin{align}
        k   & =\max\bigl\{2(k_K+4b(8+k_K)),\,2(8+4b_{\rm prot}k_K)\bigr\},   &
        C_1 & =C_\Sigma(2b,b_{\rm prot},5042\,b_{\rm prot},q,\max\{G,0\},G),
    \end{align}
    \begin{align}
        \chi        & =8(2b_{\rm prot}+1)\,b_{\rm prot}\cdot2b_{\rm prot}(b_{\rm prot}-1)
        =685013510424821760,                                                              \\
        C_{\rm net} & =1+2\cdot5042\,C_E+2C_E,                                            \\
        L_0         & =\max\bigl\{\lceil\max\{L_{\rm ang},1\}\rceil,
        \lceil2C_{\rm net}\rceil\bigr\},
    \end{align}
    \begin{align}
        c_{\rm PEPS}  & =k\chi+1,              &
        C_{\rm large} & =\max\{C_1,1\}^{\chi}.
    \end{align}
    Writing $k(c)$ and $c_{\rm PEPS}(c)=k(c)\chi+1>0$ as functions of a general patch exponent $c\in\R$, one has $c_{\rm PEPS}=c_{\rm PEPS}(c(q,J,\Delta,20D,60,C_A))$.
\end{definition}

% Source: 07-assembly.tex, lines 9--206.
% Revision: adc7f1241b42e322a6451854ab7e4b4c146bf78a.
\begin{theorem}[The polynomial PEPS approximation from angular and patch bounds]
    \label{thm:polynomial_peps_of_patch_supply}
    \lean{TNLean.PEPS.Approximation.hasPEPSApproximation_of_patchSupply}
    \leanok
    \uses{def:polynomial_peps_explicit_constants,def:peps_angular_scale,
        def:peps_adaptive_patch_constraint,def:peps_adaptive_patch_constants}
    Fix $q\ge2$ and a family of unit vectors on the $L\times L$ grids. Let $D\ge1$ and $L_{\rm ang}$ be an angular slope and threshold for the family at accuracy $L^{-60}$, clearance $a_0/2$ ($a_0=1/5000$), diameter ratio $1$, and at most $4$ marked points, and suppose every family state at side $L\ge2$ has adaptive patches on the cell-center sites at $B_0=20D$, power $60$, prefactor $C(q,20D)$, and exponent $c\in\R$. Then for every $L\ge\max\{2,L_0\}$, every family state $\Omega$ at side $L$ admits a nonzero PEPS $\Phi$ with bond dimensions $D_e\le C_{\rm large}L^{c_{\rm PEPS}(c)}$ and
    \begin{align}
        \min_{\theta\in\mathbb R}
        \left\|\frac{\Phi}{\|\Phi\|}-e^{i\theta}\Omega\right\| & \le L^{-1}.
    \end{align}
\end{theorem}

\begin{proof}\leanok
    \uses{lem:peps_dyadic_protocol_at_given_scale,
        thm:peps_large_size_network_from_protocol,
        thm:peps_ground_state_approximation_conditional}
    For $L\ge\lceil\max\{L_{\rm ang},1\}\rceil$, Lemma~\ref{lem:peps_dyadic_protocol_at_given_scale} at $C_N=5042$, $C_K=G$, and $C_E=\max\{10080,2b+1\}$ yields the protocol; Theorem~\ref{thm:peps_large_size_network_from_protocol} at accuracy $L^{-5}$ compresses it to an operator network with error $C_{\rm net}L^{-5}$, link dimensions at most $C_1L^{k(c)}$, block size at most $b_{\rm prot}$, at most $2b_{\rm prot}(b_{\rm prot}-1)$ links per party, and separation at most $b_{\rm prot}$; and for $L\ge\lceil2C_{\rm net}\rceil$, Theorem~\ref{thm:peps_ground_state_approximation_conditional} routes this network into a PEPS with $D_e\le\max\{C_1,1\}^\chi L^{k(c)\chi+1}$ and error at most $L^{-1}$.
\end{proof}

% Source: 07-assembly.tex, lines 9--206.
% Revision: adc7f1241b42e322a6451854ab7e4b4c146bf78a.
\begin{theorem}[The polynomial PEPS approximation at given area-law and angular constants]
    \label{thm:polynomial_peps_of_ratios}
    \lean{TNLean.PEPS.Approximation.polynomialPEPSApproximation_of_ratios}
    \leanok
    \uses{def:polynomial_peps_explicit_constants,def:polynomial_peps_model,
        def:square_entropy_area_law_at,def:peps_angular_scale}
    Fix $q\ge2$ and $J,\Delta>0$. Assume the square-grid entropy area law with constant $C_A$ and an angular slope $D\ge1$ and threshold $L_{\rm ang}$ (at accuracy $L^{-60}$, clearance $a_0/2$ with $a_0=1/5000$, diameter ratio $1$, and at most $4$ marked points) for the ground states of the model at $(J,\Delta)$, and let $c_{\rm PEPS},C_{\rm large},L_0$ be as in Definition~\ref{def:polynomial_peps_explicit_constants}. Then for every $L\ge\max\{2,L_0\}$ and unit ground state $\Omega$ of a model Hamiltonian (Theorem~\ref{thm:openai_polynomial_peps}) with gap $\Delta$, there is a nonzero PEPS $\Phi$ on the $L\times L$ grid with $D_e\le C_{\rm large}L^{c_{\rm PEPS}}$ and
    \begin{align}
        \min_{\theta\in\mathbb R}
        \left\|\frac{\Phi}{\|\Phi\|}-e^{i\theta}\Omega\right\| & \le L^{-1}.
    \end{align}
\end{theorem}

\begin{proof}\leanok
    \uses{thm:peps_adaptive_patch_supply_of_area_law_at,
        thm:polynomial_peps_of_patch_supply}
    Theorem~\ref{thm:peps_adaptive_patch_supply_of_area_law_at} supplies adaptive patches from $C_A$ at $B_0=20D$ and power $60$ with constants $(C,c)$ on the cell-center sites, so Theorem~\ref{thm:polynomial_peps_of_patch_supply} gives the bound with $c_{\rm PEPS}(c)=c_{\rm PEPS}$.
\end{proof}

% Status: auxiliary result. Theorem 1.1 of the source asserts only that the
% constants exist; that statement is thm:openai_polynomial_peps.
% Source: 00-introduction.tex, Theorem 1.1, and its proof in 07-assembly.tex,
% lines 9--214.
% Revision: adc7f1241b42e322a6451854ab7e4b4c146bf78a.
\begin{definition}[The constants of the polynomial PEPS approximation]
    \label{def:polynomial_peps_closed_constants}
    \lean{TNLean.PEPS.Approximation.pepsAreaLawConstant,
        TNLean.PEPS.Approximation.pepsAngularScale,
        TNLean.PEPS.Approximation.pepsAngularThreshold,
        TNLean.PEPS.Approximation.pepsSizeThreshold,
        TNLean.PEPS.Approximation.pepsBondExponent,
        TNLean.PEPS.Approximation.pepsBondPrefactor,
        TNLean.PEPS.Approximation.pepsBondExponent_pos,
        TNLean.PEPS.Approximation.pepsBondPrefactor_pos}
    \leanok
    \uses{def:polynomial_peps_explicit_constants,def:area_law_closed_constant,
        def:peps_angular_information_constants,
        def:peps_closed_information_constants}
    Fix $q\in\Z$, $J,\Delta\in\R$, and $a_0=1/5000$. Set $C_A=C_{\rm AL}(q,1,J,\Delta)$ (Definition~\ref{def:area_law_closed_constant}), $D=D_{\rm ang}(J,\Delta,a_0/2,60)$, and $L_{\rm ang}=L_{\rm ang}(q,J,\Delta,a_0/2,1,4,60)$ (Definition~\ref{def:peps_closed_information_constants}). With $L_0,c_{\rm PEPS},C_{\rm large}$ given by Definition~\ref{def:polynomial_peps_explicit_constants} at $(C_A,D,L_{\rm ang})$, define
    \begin{align}
        C_{\rm PEPS} & =\max\bigl\{C_{\rm large},\,q^{L_0^2}\bigr\}>0.
    \end{align}
\end{definition}

% Status: auxiliary result with no hypothesis beyond those of Theorem 1.1,
% whose statement is thm:openai_polynomial_peps.
% Source: 00-introduction.tex, Theorem 1.1, and its proof in 07-assembly.tex,
% lines 9--214.
% Revision: adc7f1241b42e322a6451854ab7e4b4c146bf78a.
\begin{theorem}[Polynomial PEPS approximation with explicit constants]
    \label{thm:polynomial_peps_explicit_constants}
    \lean{TNLean.PEPS.Approximation.polynomialPEPSApproximation_explicit}
    \leanok
    \uses{def:polynomial_peps_closed_constants,def:polynomial_peps_model}
    Fix $q\ge2$ and $J,\Delta>0$, and let $C_{\rm PEPS},c_{\rm PEPS}$ be the constants of Definition~\ref{def:polynomial_peps_closed_constants}. For every $L\ge2$ and unit ground state $\Omega$ of a model Hamiltonian (Theorem~\ref{thm:openai_polynomial_peps}) with gap $\Delta$, there is a nonzero PEPS $\Phi$ on the $L\times L$ grid with $D_e\le C_{\rm PEPS}L^{c_{\rm PEPS}}$ for every edge $e$ and
    \begin{align}
        \min_{\theta\in\mathbb R}
        \left\|\frac{\Phi}{\|\Phi\|}-e^{i\theta}\Omega\right\| & \le L^{-1}.
    \end{align}
\end{theorem}
\begin{proof}
    \leanok
    \uses{thm:square_area_law_closed_constant,
        thm:peps_information_from_gap_closed,thm:polynomial_peps_of_ratios,
        lem:peps_remove_small_sizes_named}
    Theorems~\ref{thm:square_area_law_closed_constant} and~\ref{thm:peps_information_from_gap_closed} establish the area-law constant $C_A$ and angular constants $D\ge1,L_{\rm ang}$, so Theorem~\ref{thm:polynomial_peps_of_ratios} gives the PEPS with prefactor $C_{\rm large}$ for $L\ge L_0$, and Lemma~\ref{lem:peps_remove_small_sizes_named} extends it to all $L\ge2$ with prefactor $C_{\rm PEPS}$.
\end{proof}
